\subsection{Ненастойчивый протокол CSMA (Non-persistent CSMA)}

В ненастойчивом протоколе CSMA терминал непрерывно слушает канал перед началом передачи. Если канал свободен, терминал немедленно начинает передачу. Если канал занят, терминал откладывает попытку на случайное время. 

Эволюция системы представляет собой чередование периодов занятости (англ.: busy period, $B$) и периодов простоя (англ.: idle period, $I$). Пара смежных интервалов образует цикл. Внутри периода $B$ между двумя последовательными передачами всегда присутствует минимум один свободный слот длины $a$. В самом деле, для инициации новой передачи терминал обязан зафиксировать свободное состояние среды.

\begin{center}
\begin{tikzpicture}[>=Stealth, scale=1.2, font=\small]
    % Ось времени
    \draw[->, thick] (-1.5,0) -- (12.5,0) node[right] {$t$};

    % Левое троеточие (предыстория системы)
    \node at (-1.0, 0.4) {$\dots$};

    % Левые два слота a (до момента начала отсчета t=0)
    \foreach \x in {-0.4, -0.2} {
        \draw[thick] (\x,0) -- (\x,0.3);
    }
    % Выровненные стрелочки для левых слотов
    \foreach \x in {-0.3, -0.1} {
        \draw[<-] (\x,0.3) -- (\x,1.2) node[above] {$a$};
    }

    % Вертикальная линия старта (t=0)
    \draw[thick] (0,-0.2) -- (0,1.8);

    % Блок 1 (Успех)
    \draw[thick, fill=SpringGreen4!10] (0,0) rectangle (1.4,0.8) node[midway] {Успех};
    \draw[thick, pattern=north east lines, pattern color=DeepSkyBlue4] (1.4,0) rectangle (1.6,0.8);
    \draw[<-] (1.5,0.8) -- (1.5,1.2) node[above] {$a$};
    % Обозначение T для Блока 1
    \draw[<->] (0,-0.12) -- (1.4,-0.12) node[midway, below, font=\scriptsize] {$T$};
    
    % Блок 2 (Коллизия)
    \draw[thick, fill=Red3!10] (1.6,0) rectangle (3.0,0.8) node[midway] {Коллизия};
    \draw[thick, pattern=north east lines, pattern color=DeepSkyBlue4] (3.0,0) rectangle (3.2,0.8);
    \draw[<-] (3.1,0.8) -- (3.1,1.2) node[above] {$a$};
    % Обозначение T для Блока 2
    \draw[<->] (1.6,-0.12) -- (3.0,-0.12) node[midway, below, font=\scriptsize] {$T$};

    % Блок 3 (Успех)
    \draw[thick, fill=SpringGreen4!10] (3.2,0) rectangle (4.6,0.8) node[midway] {Успех};
    \draw[thick, pattern=north east lines, pattern color=DeepSkyBlue4] (4.6,0) rectangle (4.8,0.8);
    \draw[<-] (4.7,0.8) -- (4.7,1.2) node[above] {$a$};
    % Обозначение T для Блока 3
    \draw[<->] (3.2,-0.12) -- (4.6,-0.12) node[midway, below, font=\scriptsize] {$T$};
    
    % Слоты простоя (I)
    \foreach \x in {4.8, 5.0, 5.2, 5.4, 5.6} {
        \draw[thick] (\x,0) -- (\x,0.3);
    }
    % Стрелочки для слотов простоя (единая высота выравнивания)
    \foreach \x in {4.9, 5.1, 5.3, 5.5} {
        \draw[<-] (\x,0.3) -- (\x,1.2) node[above] {$a$};
    }
    
    % Троеточие в середине (переход к следующему циклу)
    \node at (6.2, 0.4) {$\dots$};

    % Второй период B
    \draw[thick, fill=SpringGreen4!10] (6.8,0) rectangle (8.2,0.8) node[midway] {Успех};
    \draw[thick, pattern=north east lines, pattern color=DeepSkyBlue4] (8.2,0) rectangle (8.4,0.8);
    \draw[<-] (8.3,0.8) -- (8.3,1.2) node[above] {$a$};
    % Обозначение T для Блока 4
    \draw[<->] (6.8,-0.12) -- (8.2,-0.12) node[midway, below, font=\scriptsize] {$T$};
    
    \draw[thick, fill=SpringGreen4!10] (8.4,0) rectangle (9.8,0.8) node[midway] {Успех};
    \draw[thick, pattern=north east lines, pattern color=DeepSkyBlue4] (9.8,0) rectangle (10.0,0.8);
    \draw[<-] (9.9,0.8) -- (9.9,1.2) node[above] {$a$};
    % Обозначение T для Блока 5
    \draw[<->] (8.4,-0.12) -- (9.8,-0.12) node[midway, below, font=\scriptsize] {$T$};
    
    % Правые слоты простоя
    \foreach \x in {10.0, 10.2, 10.4, 10.6} {
        \draw[thick] (\x,0) -- (\x,0.3);
    }
    
    % Правое троеточие
    \node at (11.4, 0.4) {$\dots$};
    
    % Подписи B и I снизу
    \draw[decorate, decoration={brace, amplitude=6pt, mirror}, thick] (0, -0.45) -- (4.8, -0.45) node[midway, below=8pt] {$B$};
    \draw[decorate, decoration={brace, amplitude=6pt, mirror}, thick] (4.8, -0.45) -- (6.8, -0.45) node[midway, below=8pt] {$I$};
    \draw[decorate, decoration={brace, amplitude=6pt, mirror}, thick] (6.8, -0.45) -- (10.0, -0.45) node[midway, below=8pt] {$B$};
\end{tikzpicture}
\end{center}

\begin{statement}[Пропускная способность ненастойчивого CSMA]
    Для сети, использующей протокол ненастойчивого CSMA, зависимость нормированной пропускной способности $S$ от абсолютной интенсивности предложенной нагрузки $\lambda$ имеет вид:
    $$ S = \frac{\lambda a T e^{-\lambda a}}{T + a - T e^{-\lambda a}}, $$
    где $T$ — время передачи пакета, $a$ — длительность слота прослушивания.
\end{statement}
\begin{sproof}
    В силу теоремы восстановления пропускная способность равна отношению среднего времени успешной передачи к средней длине цикла:
    $$ S = \frac{\E[U]}{\E[B] + \E[I]}, $$
    где $\E[B]$ и $\E[I]$ — математические ожидания длительностей периодов занятости и простоя соответственно, а $\E[U]$ — среднее суммарное время успешных передач за один цикл.

    Вычислим математическое ожидание периода простоя $\E[I]$. Период простоя состоит из дискретного числа слотов длины $a$. Он включает ровно $k$ слотов, если в течение первых $k-1$ слотов заявки не поступали, а в $k$-м слоте появилась хотя бы одна заявка. Так как входящий поток пуассоновский с интенсивностью $\lambda$, вероятность отсутствия заявок на интервале $a$ равна $e^{-\lambda a}$. Следовательно, число слотов подчиняется геометрическому распределению:
    $$ \P(I = ka) = \left(e^{-\lambda a}\right)^{k-1} (1 - e^{-\lambda a}), \quad k \ge 1. $$
    Отсюда математическое ожидание длины периода простоя составляет:
    $$ \E[I] = \frac{a}{1 - e^{-\lambda a}}. $$

    Рассмотрим период занятости $B$. Он представляет собой последовательность тактов передачи (англ.: transmission period, $TP$). Каждый такт включает непосредственную передачу пакета (длительность $T$) и обязательный слот прослушивания (длительность $a$), откуда полная длина такта равна $T+a$. Новый такт инициируется, если в течение слота $a$ в систему поступает пакет. В противном случае начинается период простоя $I$. Вероятность того, что период занятости состоит из $k$ тактов, равна:
    $$ \P(B = k(T + a)) = \left(1 - e^{-\lambda a}\right)^{k-1} e^{-\lambda a}, \quad k \ge 1. $$
    Значит, математическое ожидание периода занятости равно:
    $$ \E[B] = \frac{T+a}{e^{-\lambda a}}. $$

    Вычислим $\E[U]$. Среднее число тактов передачи в одном периоде $B$ составляет $\E[B] / (T+a) = 1/e^{-\lambda a}$. Передача завершается успехом тогда и только тогда, когда в предшествующем слоте $a$ поступила ровно одна заявка. Поскольку такт передачи уже состоялся, достоверно известно, что число заявок не равно нулю. Следовательно, вероятность успеха $P_{suc}$ вычисляется как условная вероятность:
    $$ P_{suc} = \P(\text{поступила 1 заявка} \mid \text{поступило } \ge 1 \text{ заявок}) = \frac{\lambda a e^{-\lambda a}}{1 - e^{-\lambda a}}. $$
    Умножая среднее число тактов на вероятность успеха и длительность полезной передачи $T$, получаем:
    $$ \E[U] = \frac{1}{e^{-\lambda a}} \cdot \frac{\lambda a e^{-\lambda a}}{1 - e^{-\lambda a}} \cdot T = \frac{\lambda a T}{1 - e^{-\lambda a}}. $$

    Подставим найденные величины $\E[I]$, $\E[B]$ и $\E[U]$ в исходную формулу для пропускной способности:
    $$ S = \frac{\frac{\lambda a T}{1 - e^{-\lambda a}}}{\frac{T+a}{e^{-\lambda a}} + \frac{a}{1 - e^{-\lambda a}}}. $$
    Домножим числитель и знаменатель дроби на величину $e^{-\lambda a} (1 - e^{-\lambda a})$. Тогда числитель преобразуется к виду $\lambda a T e^{-\lambda a}$, а знаменатель принимает вид:
    $$ (T+a)(1 - e^{-\lambda a}) + a e^{-\lambda a} = T + a - T e^{-\lambda a} - a e^{-\lambda a} + a e^{-\lambda a} = T + a - T e^{-\lambda a}. $$
    Окончательно получаем требуемое соотношение для $S$.
\end{sproof}

\begin{conseq}
    Перейдем к пределу при устремлении длины слота прослушивания к нулю ($a \to 0$). Разделим числитель и знаменатель полученной дроби на $a$:
    $$ S = \frac{\lambda T e^{-\lambda a}}{\frac{T}{a}(1 - e^{-\lambda a}) + 1}. $$
    Используя первый замечательный предел для экспоненты $\lim\limits_{a \to 0} \frac{1 - e^{-\lambda a}}{a} = \lambda$, находим асимптотическую пропускную способность идеализированной системы:
    $$ \lim_{a \to 0} S = \frac{\lambda T}{\lambda T + 1} = \frac{G}{G + 1}, $$
    где $G = \lambda T$ — среднее число заявок, генерируемых за время передачи одного пакета. При неограниченном росте предложенной нагрузки ($G \to \infty$) предел функции равен $1$. Это показывает, что теоретически ненастойчивый CSMA способен обеспечить абсолютную утилизацию канала без потерь на коллизии.
\end{conseq}

\begin{center}
    \begin{tikzpicture}
        \begin{axis}[
            width=0.8\textwidth, height=8cm,
            xlabel={Предложенная нагрузка $G = \lambda T$}, 
            ylabel={Пропускная способность $S$},
            xmin=0, xmax=10, ymin=0, ymax=1.1,
            grid=major,
            samples=150,
            domain=0.01:10,
            legend pos=south east,
            legend style={font=\small},
            thick,
            title={Зависимость пропускной способности от нагрузки (при $T=1$)}
        ]
            % a = 0
            \addplot[black, dashed] {x/(1+x)}; 
            \addlegendentry{$a=0$}
            
            % a = 0.1T
            \addplot[Red3] {x*0.1*exp(-x*0.1)/(1 - exp(-x*0.1) + 0.1)}; 
            \addlegendentry{$a=0.1T$}
            
            % a = 0.2T
            \addplot[Tan1] {x*0.2*exp(-x*0.2)/(1 - exp(-x*0.2) + 0.2)}; 
            \addlegendentry{$a=0.2T$}

            % a = 0.5T
            \addplot[DeepSkyBlue4] {x*0.5*exp(-x*0.5)/(1 - exp(-x*0.5) + 0.5)}; 
            \addlegendentry{$a=0.5T$}
            
            % a = 1.0T
            \addplot[SpringGreen4] {x*1.0*exp(-x*1.0)/(1 - exp(-x*1.0) + 1.0)}; 
            \addlegendentry{$a=T$}
        \end{axis}
    \end{tikzpicture}
\end{center}